\chapter{OOP and UML}
\label{chap:c++newfeatures}
\vspace{10.5cm}
%%% OOP.
\section{Object Oriented Programming}
    \subsection{Polymorphism}
    In C++ there are different kinds of polymorphism, summarized it 
    the figure~\ref{fig:cpp-poly}
    \begin{figure}[H]
        \centering
        \includegraphics[scale=0.5]{OOPandUML/cpp_polymorphism}
        \caption{Polimorphism in C++}
        \label{fig:cpp-poly}
    \end{figure}
    \subsubsection*{Compile time polymorphism}
    Overloaded functions are invoked by matching the type and number of 
    arguments at compile time, because the compiler already has all the 
    information needed. This is know by the name of 
    \textbf{static binding} or \textbf{}.
    \subsubsection*{Run time polymorphism}
%%% UML 
\section{UML}
    %% Class Diagram
    \subsection{Class Diagram}
    In Unified Modeling Language, a class diagram is used to describe 
    the static structure of a system by showing the system's classes, 
    their attributes, operations (methods) and the relationship among 
    them.
    \newline
    The purpose of a class diagram is the conceptual modeling of the 
    structure of the software application. In case of detailed modeling, 
    it is also possible to translate the diagram into programming code.
    %%% Visibility
    \subsubsection{Visibility}
    To indicate visibility for methods and sttributes, this notation 
    must be placed before their names
    \begin{table}[ht]
        \begin{tabular}{l l}
            $+$ \quad \quad & public\\
            $-$ & private\\
            \texttt{\#} & protected\\
            \textasciitilde & package\\
        \end{tabular}
    \end{table}
    %%% Scope
    \subsubsection{Scope}
    UML defines \textbf{instance} and \textbf{class} scopes for members. 
    The latter is represented by \underline{underlined names}.
    \begin{itemize}
        \item \textbf{instance members} are scoped to a specific instance
        \item \textbf{class members} are the one commonly recognized as 
        "static"
    \end{itemize}
    %%% Arrows
    \subsubsection{Arrows}
    The relationship defined in UML are represented by the following 
    logical connectors
    \begin{figure}[H]
        \centering
        \includegraphics[scale=0.5]{OOPandUML/uml_arrow_cheatsheet}
        \caption{UML logical connectors}
        \label{fig:uml-arrows}
    \end{figure}
    \begin{itemize}
        \item \textbf{Dependency} is a relationship which shows that an 
              element requires other model elements for their 
              specification or implementation. 
              \newline
              The dashed lines points from the dependent (client) to 
              the independent (supplier) element, specifying the 
              direction of a relationship and not of a process.
        \item \textbf{Association} defines a relationship bewteen classes 
              of objects that allows one object instance to cause another 
              to perform an action on its behalf. This relationship is 
              structural and does not represent any behavioural trait.
              \newline
              Code implementation usually requires the requsting object 
              to invoke a method or member function using a 
              \textbf{reference} or a \textbf{pointer} to the controlled 
              object.
              \newline
              Association can be \textbf{bi}-directional, 
              \textbf{uni}-directional, aggregation and reflexive.
              \newline
              Objects related via the association are considered to act 
              in a role with respect to the association.
              \newline{}
              The ends of the association can have all the 
              characteristics of a property: multiplicity, name, 
              visibility and a specification explaining whether the 
              association is ordered and / or unique.
        \item \textbf{Aggregation} is a variant of "\textbf{has a}" 
              association relationship and more specific than the latter. 
              It represents a part-whole or a part-of relationship. 
              As a type of association, an aggregation can be named and 
              have the same adornments an association can.
              \newline
              An aggregation can occur when a class is a collection or 
              a container of other classes, but the contained classed 
              do not have a strong lifecycle dependency on the 
              container. \textbf{The content of the container continues 
              to exist when the container is destroyed}.
              \newline
              In UML notation, the hollow diamond shape is drawn on the 
              containing class end of the connecting line. An example of 
              aggregarion is a Pond which has zero or more Ducks 
              and Duck which has at most one Pond. Duck can exist 
              separately from a pond.
        \item \textbf{Composition} refers to combining objects within 
              compound objects, ensuring the encapsulation of each 
              object by using the well-defined interface without 
              exposing the internals. Differently, data structures do 
              not enforce encapsulation.
              \newline
              Composition can be regarded as a relationship between 
              types: an object of composite type "\textbf{has}" objects 
              of other types
              \newline
              Composition must be distinguished from subtyping, which 
              is the process of adding detail to a general data type 
              to create a more specific one.
              \newline
              In UML notation, the filled diamond shape is drawn on the 
              containing class end of the connecting line. An example of 
              composition is a Univerity which has one or more 
              Departments, which exactly one Univerity they belong 
              to. A Department can not exist as a separate entity 
              detached from a specific University.
    \end{itemize}
    %%% Composition versus Aggregation
    %~ \nparagraph{Composition versus Aggregation}
    %~ \begin{itemize}
        %~ \item Composition relationship
        %~ \item Aggregation relationship
    %~ \end{itemize}
    %%% Class level relationship
    \subsubsection{Class-level relationships}
    \begin{itemize}
        \item \textbf{Generalization/Inheritance} indicates the one of 
              the two related classes (called \textbf{subclass}) is 
              considered to be a specialization from the other (called 
              \textbf{super type}). Therefore any instance of the 
              subtype is also an instance of the superclass. An example 
              comes from biological classification: humans are a 
              subclass of simian which is a subclass of mammal.
              \newline
              Generalization is also known as \textbf{inheritance} or a 
              "\textbf{is a}" relationship. 
              \newline
              The base class is also known as parent, superclass, 
              base class or base type.
              \newline
              The subtype is also known as child, subclass, 
              derived class, derived type, inheriting class or 
              inheriting type.
        \item \textbf{Realization/Implementation} is a relationship 
              between two model elements: one model element (called 
              client) realizes (implements or executes) the behavior 
              that the other model element (the supplier) specifies.
    \end{itemize}
    %%% General relationship
    \subsubsection{General relationship}
    \begin{itemize}
        \item \textbf{Dependency} can be a weaker form of bond that 
              indicates that one class depends on another because it 
              uses it at some point in time. Sometimes this relationship 
              might be implemented as member function arguments.
        \item \textbf{Multiplicity} is an optional notation which can be 
              useful to add details to the association relationship. At 
              each end of the line, one of the following notation can be 
              used to indicate the range of number of objexts that 
              partecipate in the association from the perspective of the 
              other end.
             \begin{table}[ht]
                 \begin{center}
                     \begin{tabular}{l l}
                         $0$ & No instances\\
                         $0..1$ \quad & No instances, or one instance\\
                         $1$ & Exactly one instance\\
                         $1..1$ & Exactly one instance\\
                         $0..*$ & Zero or more instances\\
                         $*$ & Zero or more instances\\
                         $1..*$ & One or more instances\\
                     \end{tabular}
                 \end{center}
             \end{table}
    \end{itemize}

    \subsubsection{Class and Sequence diagram Cheatsheet}
    Following a summary of the information for class and sequance UML 
    diagrams
    \begin{figure}[H]
        \centering
        \includegraphics[scale=0.7]{OOPandUML/uml_class_sequence_diag_cheatsheet}
        \caption{Class and Sequence diagrams cheatsheet}
        \label{fig:class-seq-diag-cheatsheet}
    \end{figure}
